\documentclass[compress, 8pt, aspectratio=169]{beamer}

\usepackage{presentationtemplate}
\usepackage[fontsize=\small,linenosfontsize=\tiny,askip=3mm]
  {mylisting}
\usepackage[fontsize=\small,askip=3mm]{terminal}
\usepackage{textcomp}
\usepackage[autostyle=true]{csquotes}

\usetikzlibrary{calc}

\title{CMake}

\newcommand{\cmakehelp}%
  {https://cmake.org/cmake/help/latest}

\newcommand{\stdfl}{-std=gnu99 -Wall -Wextra}

\newcommand{\dmtc}{toolchain/demo-toolchain.cmake}

\newcommand{\armtc}{toolchain/aarch64-toolchain.cmake}

\makeatletter
\newdimen\cm@FVleftmargin
\let\cm@FVSetupFont\FV@SetupFont
\def\FV@SetupFont{%
  \cm@FVleftmargin=\leftmargin
  \cm@FVSetupFont
  \let\leftmargin\cm@FVleftmargin}
\makeatother

\begin{document}

\begin{frame}
  \titlepage
\end{frame}

\begin{frame}{Проблемы GNU Make}

  \begin{itemize}

    \item Большинство типичных действий (поиск библиотек, добавление
      флагов компиляции к уже существующим, проверка работоспособности
      компилятора) требует написания кода.

    \item Трудно изменять директорию сборки.

    \item Трудно поддерживать несколько компиляторов и платформ сборки.

    \item Не очень хорошо поддерживается модульность: функциональности
      подкаталогов иногда недостаточно.

    \item Управление зависимостями реализовать трудно.

  \end{itemize}

  \bigbreak

  GNU Make не очень хорошо масштабируется на большие и сложные проекты.

\end{frame}

\begin{frame}{CMake как мета-система сборки}

  \begin{itemize}

    \item CMake описывает проект в \texttt{CMakeLists.txt} на собственном языке.

    \item CMake описывает сборку проекта в виде набора \textit{целей сборки}.
      Целями являются артефакты: исполняемые файлы, библиотеки и т.д.

    \item CMake \textit{генерирует} файлы сборки для выбранного сборочного
      бэкенда (генератора): \texttt{Makefile}, \texttt{build.ninja}, проекты
      Visual Studio/Xcode\footnotemark{}.

      \vspace{0.6em}

      \begin{center}

        \begin{tikzpicture}[
          box/.style={draw, thick, align=center, inner sep=4pt,
            font=\footnotesize}
        ]

          \node[box] (src) {\texttt{CMakeLists.txt}};

          \node[box, right=2.4cm of src] (gen)
            {\texttt{Makefile} / \texttt{build.ninja} / \texttt{*.vcxproj}};

          \draw[->, thick] (src.east) -- node[above, font=\footnotesize]
            {генерация} (gen.west);

        \end{tikzpicture}

      \end{center}

      \vspace{0.6em}

    \item CMake более \textit{декларативен}, чем типичная система сборки
      для C/C++.

  \end{itemize}

  \footnotetext{\url{\cmakehelp/manual/cmake.1.html}}

\end{frame}

\begin{frame}{CMake и GNU Make: сравнение}

  \begin{center}

    \renewcommand{\arraystretch}{1.4}

    \begin{tabular}{|p{2.6cm}|p{5.0cm}|p{5.0cm}|}

      \hline

      \textbf{Критерий} & \textbf{GNU Make} & \textbf{CMake} \\

      \hline

      Уровень описания
        & рецепт: \textbf{как} (команды, флаги)
        & цель: \textbf{что} (библиотеки, исполняемый файлы) \\

      \hline

      Установка тулчейна
        & вручную \texttt{make CC=\dots}
        & автоматически или через toolchain-файл \\

      \hline

      Кто пишет команды компиляции
        & разработчик
        & сгенерированы \\

      \hline

    \end{tabular}

  \end{center}

\end{frame}

\begin{frame}{Этапы сборки CMake-проекта}

  \begin{enumerate}

    \item \textbf{configure} --- CMake находит компилятор, применяет toolchain-файл
      (если задан), обрабатывает опции сборки, создает кэш и
      генерирует файлы сборки.

    \item \textbf{build} --- CMake делегирует сборочному бэкенду выполнение
      непосредственно компиляции.

    \item \textbf{install} --- CMake делегирует сборочному бэкенду установку.

  \end{enumerate}

\end{frame}

\begin{frame}{Этапы сборки CMake-проекта}

  \begin{center}

    \begin{tikzpicture}[
      box/.style={draw, thick, align=center, inner sep=4pt,
        font=\footnotesize}
    ]

      \node[box] (src) {исходники + \texttt{CMakeLists.txt}};

      \node[box, below=1.0cm of src] (bld)
        {\texttt{build/}: \texttt{CMakeCache.txt} +
        \texttt{Makefile}/\texttt{build.ninja}};

      \node[box, below=1.0cm of bld] (art)
        {артефакты: \texttt{libcalc.a}, \texttt{calc\_app}};

      \node[box, below=1.0cm of art] (pre)
        {установка в: \texttt{bin/}, \texttt{lib/}, \texttt{include/}};

      \draw[->, thick] (src.south) -- node[right, align=left,
        font=\footnotesize] {1. \textbf{configure}: \texttt{cmake -S . -B build/}}
        (bld.north);

      \draw[->, thick] (bld.south) -- node[right, align=left,
        font=\footnotesize] {2. \textbf{build}: \texttt{cmake --build build/}}
        (art.north);

      \draw[->, thick] (art.south) -- node[right, align=left,
        font=\footnotesize] {3. \textbf{install}: \texttt{cmake --install build/}}
        (pre.north);

    \end{tikzpicture}

  \end{center}

\end{frame}

\begin{frame}[fragile]{Синтаксис \texttt{CMakeLists.txt}}

  \myinputlisting[minted language=cmake]
    {presentations/02-cmake/src/}
    {syntax-demo/CMakeLists.txt}

  \bigbreak

  \begin{terminalwindow}
!\shellcommand{cmake -S syntax-demo/ -B \$(mktemp -d)}!
-- The C compiler identification is GNU 15.3.0
-- Detecting C compiler ABI info
-- Detecting C compiler ABI info - done
-- Check for working C compiler: !\dots!/bin/gcc - skipped
-- Detecting C compile features
-- Detecting C compile features - done
-- flag: -O2
-- Configuring done (0.9s)
-- Generating done (0.0s)
-- Build files have been written to: /tmp/!\dots!
  \end{terminalwindow}

\end{frame}

\begin{frame}[fragile]{Простейший \texttt{CMakeLists.txt}}

  \begin{columns}[T]

    \begin{column}{0.40\textwidth}

      \myinputlisting[minted language=cmake]
        {presentations/02-cmake/src/}
        {calc/01/CMakeLists.txt}

    \end{column}

    \begin{column}{0.56\textwidth}

      \begin{terminalwindow}
!\shellcommand{cmake -S calc/01/ -B /tmp/calc-01}!
-- The C compiler identification is GNU 15.3.0
-- Detecting C compiler ABI info
-- Detecting C compiler ABI info - done
-- Check for working C compiler: !\dots!/bin/gcc - skipped
-- Detecting C compile features
-- Detecting C compile features - done
-- Configuring done (0.9s)
-- Generating done (0.0s)
-- Build files have been written to: /tmp/calc-01
!\shellcommand{cmake --build /tmp/calc-01/}!
[ 25%] Building C object CMakeFiles/calc.dir/calc.c.o
[ 50%] Linking C static library libcalc.a
[ 50%] Built target calc
[ 75%] Building C object CMakeFiles/calc_app.dir/main.c.o
[100%] Linking C executable calc_app
[100%] Built target calc_app
      \end{terminalwindow}

    \end{column}

  \end{columns}

\end{frame}

\begin{frame}[fragile]{Простейший \texttt{CMakeLists.txt}}

  \begin{columns}[T]

    \begin{column}{0.48\textwidth}

      \myinputlisting[minted language=c]
        {presentations/02-cmake/src/}
        {calc/01/calc.h}

    \end{column}

    \begin{column}{0.48\textwidth}

      \myinputlisting[minted language=c]
        {presentations/02-cmake/src/}
        {calc/01/calc.c}

      \vspace{0.5em}

      \myinputlisting[minted language=c]
        {presentations/02-cmake/src/}
        {calc/01/main.c}

    \end{column}

  \end{columns}

\end{frame}

\begin{frame}[fragile]{Цели сборки}

  \begin{itemize}

    \item \textbf{Цель (target)} --- именованный узел графа
      зависимостей: библиотека или исполняемый файл, с набором свойств
      (включённые каталоги, флаги, подключаемые библиотеки, имя
      выходного файла).

    \item Все настройки привязываются к \textbf{цели} через
      \texttt{target\_*}-команды\footnotemark{}.

    \item Имя цели $\ne$ имя файла: \texttt{calc} $\rightarrow$
      \texttt{libcalc.a}; \texttt{calc\_app} $\rightarrow$
      \texttt{calc\_app} (имя файла выводится из имени цели и типа).

  \end{itemize}

  \vspace{0.6em}

  \begin{center}

    \begin{tikzpicture}[
      box/.style={draw, thick, align=center, inner sep=4pt,
        font=\footnotesize}
    ]

      \node[box] (app) {\texttt{calc\_app}};

      \node[box, right=3.0cm of app] (lib) {\texttt{calc}};

      \draw[->, thick] (app.east) -- node[above, font=\footnotesize]
        {\texttt{линкуется с}} (lib.west);

    \end{tikzpicture}

  \end{center}

  \footnotetext{\url{\cmakehelp/manual/cmake-buildsystem.7.html}}

\end{frame}

\begin{frame}[fragile]{\texttt{PUBLIC} / \texttt{PRIVATE} /
  \texttt{INTERFACE}}

  \texttt{scope} в \texttt{<command>(<target> <scope> \dots)} определяет, как
  используются ли настройки цели \texttt{target} (include directories, libraries
  и т.д.) ее потребителями.

  \vspace{0.5em}

  \begin{center}

    \renewcommand{\arraystretch}{1.4}

    \begin{tabular}{|p{3.2cm}|p{4.7cm}|p{4.7cm}|}

      \hline

      \textbf{Scope} & \textbf{Цель}
        & \textbf{Потребитель} \\

      \hline

      \texttt{PRIVATE} & да & нет \\

      \hline

      \texttt{PUBLIC} & да & да \\

      \hline

      \texttt{INTERFACE} & нет & да \\

      \hline

    \end{tabular}

  \end{center}

  \bigbreak

  \begin{myinplacelisting}[minted language=cmake]
add_library(A a.c)

target(B b.c)
target_include_libraries(A PUBLIC B)

target_include_directories(A PUBLIC ./headers) # both A and B get -I ./headers
target_compile_options(A INTERFACE -dummy) # only B gets -dummy flag
target_link_libraries(A PRIVATE foo) # only A gets -Lfoo
  \end{myinplacelisting}

\end{frame}

\begin{frame}[fragile]{Динамическая библиотека: \texttt{SHARED}}

  \begin{columns}[T]

    \begin{column}{0.52\textwidth}

      \myinputlisting[minted language=cmake]
        {presentations/02-cmake/src/}
        {calc/05/CMakeLists.txt}

    \end{column}

    \begin{column}{0.46\textwidth}

      \begin{itemize}

        \item Меняем одно слово:
          \texttt{add\_library(calc SHARED calc.c)}\footnotemark{} ---
          остальное то же (листинг
          \texttt{calc/05/CMakeLists.txt}).

        \item \texttt{-fPIC} и \texttt{-shared} CMake добавляет
          \textbf{сам} --- в лекции 1 мы писали их руками в Makefile.

        \item Артефакт --- \texttt{libcalc.so}; исполняемый файл
          хранит лишь имя --- загрузка в момент запуска.

      \end{itemize}

    \end{column}

  \end{columns}

  \bigbreak

  \begin{terminalwindow}
!\shellcommand{ldd build-so/calc\_app | grep libcalc}!
        libcalc.so => /tmp/opencode/cmake-calc/build-so/libcalc.so (!\dots!)
!\shellcommand{readelf -d build-so/calc\_app | grep RUNPATH}!
  0x000000000000001d (RUNPATH)   Library runpath: [/tmp/!\dots!/build-so:!\dots!]
  \end{terminalwindow}

  \footnotetext{\url{\cmakehelp/command/add_library.html}}

\end{frame}

\begin{frame}[fragile]{Динамическая библиотека: \texttt{SHARED}}

  Из build-дерева \texttt{./build/calc\_app} работает сразу: CMake
  прописывает RUNPATH на каталог сборки (\texttt{readelf}); после
  установки --- системные каталоги или rpath (как в лекции 1).

  \bigbreak

  \begin{terminalwindow}
!\shellcommand{cmake -S . -B build-so \&\& cmake --build build-so}!
[ 25%] Building C object CMakeFiles/calc.dir/calc.c.o
[ 50%] Linking C shared library libcalc.so
[ 50%] Built target calc
[ 75%] Building C object CMakeFiles/calc_app.dir/main.c.o
[100%] Linking C executable calc_app
[100%] Built target calc_app
!\shellcommand{ldd build-so/calc\_app | grep libcalc}!
        libcalc.so => /tmp/opencode/cmake-calc/build-so/libcalc.so (!\dots!)
!\shellcommand{readelf -d build-so/calc\_app | grep RUNPATH}!
  0x000000000000001d (RUNPATH)   Library runpath: [/tmp/!\dots!/build-so:!\dots!]
  \end{terminalwindow}

\end{frame}

\begin{frame}[fragile]{Флаги и стандарт языка}

  \myinputlisting[minted language=cmake]
    {presentations/02-cmake/src/}
    {calc/02/CMakeLists.txt}

  \bigbreak

  Что видно в реальной командной строке (проверено
  \texttt{cmake --build --verbose}):

  \vspace{0.5em}

  \begin{terminalwindow}
cc -I!\dots!/calc -O3 -DNDEBUG !\myhighlight{yellow}{\stdfl}! -MD -o CMakeFiles/!\dots! -c calc.c
cc -O3 -DNDEBUG CMakeFiles/calc_app.dir/main.c.o -o calc_app libcalc.a
  \end{terminalwindow}

\end{frame}

\begin{frame}[fragile]{\texttt{option()} и \texttt{if()}}

  Один \texttt{CMakeLists.txt} --- две сборки: \texttt{option()} задаёт
  кэш-переменную с дефолтом, \texttt{if()} выбирает
  \texttt{SHARED}/\texttt{STATIC}.

  \vspace{0.5em}

  \myinputlisting[minted language=cmake]
    {presentations/02-cmake/src/}
    {calc/03/CMakeLists.txt}

\end{frame}

\begin{frame}[fragile]{\texttt{option()} и \texttt{if()}}

  \begin{terminalwindow}
!\shellcommand{cmake -S . -B build-shared -DBUILD\_SHARED=ON}!
!\shellcommand{cmake --build build-shared}!
!\dots!
[ 50%] Linking C shared library libcalc.so
[100%] Built target calc_app
!\shellcommand{./build-shared/calc\_app}!
2 + 3 = 5
!\dots!
  \end{terminalwindow}

\end{frame}

\begin{frame}[fragile]{Подкаталоги: \texttt{add\_subdirectory()}}

  Проект распадается по каталогам, CMake делает это сам ---
  через \texttt{add\_subdirectory()}:

  \vspace{0.5em}

  \myinputlisting[minted language=text]
    {presentations/02-cmake/src/}
    {calc-subdir-tree.txt}

\end{frame}

\begin{frame}[fragile]{Подкаталоги: \texttt{add\_subdirectory()}}

  \begin{columns}[T]

    \begin{column}{0.48\textwidth}

      \myinputlisting[minted language=cmake]
        {presentations/02-cmake/src/}
        {calc-subdir/CMakeLists.txt}

    \end{column}

    \begin{column}{0.48\textwidth}

      \myinputlisting[minted language=cmake]
        {presentations/02-cmake/src/}
        {calc-subdir/lib/CMakeLists.txt}

    \end{column}

  \end{columns}

  \bigbreak

  \begin{terminalwindow}
!\shellcommand{cmake -S calc-subdir -B build-subdir -DCMAKE\_BUILD\_TYPE=Release}!
!\dots!
!\shellcommand{cmake --build build-subdir}!
[ 50%] Linking C static library libcalc.a
[100%] Linking C executable calc_app
!\shellcommand{./build-subdir/app/calc\_app}!
2 + 3 = 5
!\dots!
  \end{terminalwindow}

\end{frame}

\begin{frame}[fragile]{Конфигурация: \texttt{cmake -S . -B build}}

  Форма: \texttt{cmake -S <source-dir> -B <build-dir>}
  (CMake $\geq$ 3.13) --- исходный и сборочный каталоги явно; каталог
  сборки создаётся сам.

  \vspace{0.5em}

  \begin{terminalwindow}
!\shellcommand{cmake -S . -B build -DCMAKE\_BUILD\_TYPE=Release}!
-- The C compiler identification is Clang 21.1.8
-- Detecting C compiler ABI info
-- Detecting C compiler ABI info - done
-- Check for working C compiler: /run/!\dots!/bin/cc - skipped
-- Detecting C compile features
-- Detecting C compile features - done
-- Configuring done (0.3s)
-- Generating done (0.0s)
-- Build files have been written to: /tmp/!\dots!/build
  \end{terminalwindow}

\end{frame}

\begin{frame}[fragile]{Что сгенерировал CMake}

  \begin{terminalwindow}
!\shellcommand{ls build}!
CMakeCache.txt  CMakeFiles  Makefile  cmake_install.cmake
  \end{terminalwindow}

  \bigbreak

  \begin{terminalwindow}
!\shellcommand{grep -E '\char"5E(CMAKE\_BUILD\_TYPE|CMAKE\_C\_COMPILER|\dots):' CMakeCache.txt}!
CMAKE_BUILD_TYPE:STRING=Release
CMAKE_C_COMPILER:FILEPATH=/run/current-system/sw/bin/cc
CMAKE_INSTALL_PREFIX:PATH=/usr/local
CMAKE_GENERATOR:INTERNAL=Unix Makefiles
  \end{terminalwindow}

\end{frame}

\begin{frame}[fragile]{Сборка: \texttt{cmake --build build}
  (инкрементальная)}

  \begin{terminalwindow}
!\shellcommand{cmake --build build}!
[ 25%] Building C object CMakeFiles/calc.dir/calc.c.o
[ 50%] Linking C static library libcalc.a
[ 50%] Built target calc
[ 75%] Building C object CMakeFiles/calc_app.dir/main.c.o
[100%] Linking C executable calc_app
[100%] Built target calc_app
!\shellcommand{cmake --build build}!
[ 50%] Built target calc
[100%] Built target calc_app
  \end{terminalwindow}

\end{frame}

\begin{frame}[fragile]{Сборка: \texttt{cmake --build build}
  (инкрементальная)}

  \begin{terminalwindow}
!\shellcommand{touch calc.c}!
!\shellcommand{cmake --build build}!
[ 25%] Building C object CMakeFiles/calc.dir/calc.c.o
[ 50%] Linking C static library libcalc.a
[ 50%] Built target calc
[ 75%] Linking C executable calc_app
[100%] Built target calc_app
!\shellcommand{touch calc.h}!
!\shellcommand{cmake --build build}!
[ 25%] Building C object CMakeFiles/calc.dir/calc.c.o
[ 50%] Linking C static library libcalc.a
[ 50%] Built target calc
[ 75%] Building C object CMakeFiles/calc_app.dir/main.c.o
[100%] Linking C executable calc_app
[100%] Built target calc_app
  \end{terminalwindow}

\end{frame}

\begin{frame}[fragile]{Генераторы: \texttt{-G}}

  Генератор --- \enquote{формат} файлов сборки, которые пишет CMake.
  На Linux по умолчанию --- Unix Makefiles; другой генератор --- флаг
  \texttt{-G}:

  \begin{center}

    \renewcommand{\arraystretch}{1.4}

    \begin{tabular}{|p{2.2cm}|p{2.7cm}|}

      \hline

      \textbf{Генератор} & \textbf{Файлы сборки} \\

      \hline

      Unix Makefiles (дефолт на Linux) & \texttt{Makefile} \\

      \hline

      Ninja & \texttt{build.ninja} \\

      \hline

      Visual Studio 17 2022 & \texttt{.sln} + \texttt{.vcxproj} \\

      \hline

      Xcode & \texttt{.xcodeproj} \\

      \hline

    \end{tabular}

  \end{center}

\end{frame}

\begin{frame}[fragile]{Генераторы: \texttt{-G}}

  \begin{terminalwindow}
!\shellcommand{cmake -S . -B build-ninja -G Ninja}!
!\dots!
-- Build files have been written to: /tmp/!\dots!/build-ninja
!\shellcommand{ls build-ninja}!
build.ninja
CMakeCache.txt
CMakeFiles
cmake_install.cmake
!\shellcommand{cmake --build build-ninja}!
[1/4] Building C object CMakeFiles/calc.dir/calc.c.o
[2/4] Building C object CMakeFiles/calc_app.dir/main.c.o
[3/4] Linking C static library libcalc.a
[4/4] Linking C executable calc_app
!\shellcommand{./build-ninja/calc\_app}!
2 + 3 = 5
  \end{terminalwindow}

\end{frame}

\begin{frame}[fragile]{Переконфигурация: один каталог --- новые опции}

  \begin{terminalwindow}
!\shellcommand{cmake -S . -B build-shared}!
-- Build files have been written to: /tmp/!\dots!/build-shared
!\shellcommand{cmake -L build-shared | grep BUILD\_SHARED}!
!\myhighlight{yellow}{BUILD\_SHARED:BOOL=OFF}!
!\shellcommand{cmake -S . -B build-shared -DBUILD\_SHARED=ON}!
-- Build files have been written to: /tmp/!\dots!/build-shared
!\shellcommand{cmake -L build-shared | grep BUILD\_SHARED}!
!\myhighlight{yellow}{BUILD\_SHARED:BOOL=ON}!
!\shellcommand{cmake --build build-shared}!
[ 50%] Linking C shared library libcalc.so
[ 50%] Built target calc
[100%] Linking C executable calc_app
[100%] Built target calc_app
!\shellcommand{ldd build-shared/calc\_app | grep libcalc}!
        libcalc.so => !\dots!/build-shared/libcalc.so
  \end{terminalwindow}

\end{frame}

\begin{frame}[fragile]{Переменные и области: \texttt{set()} и
  \texttt{CACHE}}

  Два \enquote{мира} переменных: обычная (область каталога) и
  кэш-переменная (\texttt{CMakeCache.txt}):

  \vspace{0.5em}

  \myinputlisting[minted language=cmake]
    {presentations/02-cmake/src/vars-demo/}
    {CMakeLists.txt}

\end{frame}

\begin{frame}[fragile]{Переменные и области: \texttt{set()} и
  \texttt{CACHE}}

  \begin{terminalwindow}
!\shellcommand{cmake -S vars-demo -B build-vars}!
-- WHO   = directory variable
-- WHERE = default
!\shellcommand{cmake -S vars-demo -B build-vars2 -DWHERE=from-cli}!
-- WHO   = directory variable
-- WHERE = from-cli
!\shellcommand{grep WHERE build-vars2/CMakeCache.txt}!
WHERE:STRING=from-cli
  \end{terminalwindow}

\end{frame}

\begin{frame}[fragile]{\texttt{install()}: что ставить и куда}

  \myinputlisting[minted language=cmake,
    minted options={breaklines, breakanywhere, linenos,
      escapeinside=||, numbersep=2pt, highlightlines={12-14}}]
    {presentations/02-cmake/src/}
    {calc/04/CMakeLists.txt}

\end{frame}

\begin{frame}[fragile]{Toolchain-файл: что это и когда читается}

  Toolchain-файл --- CMake-скрипт: целевая платформа, компиляторы,
  sysroot; передаётся переменной \texttt{CMAKE\_TOOLCHAIN\_FILE}:

  \vspace{0.5em}

  \begin{center}

    \begin{tikzpicture}[
      box/.style={draw, thick, align=center, inner sep=3pt,
        font=\footnotesize, minimum height=0.7cm},
      arr/.style={->, thick}
    ]

      \node[box] (a) {\texttt{CMakeLists.txt}};

      \node[box, below=0.6cm of a] (b)
        {\texttt{project()}};

      \node[box, below=0.6cm of b] (c)
        {читается toolchain-файл};

      \node[box, below=0.6cm of c] (d)
        {определение компиляторов};

      \draw[arr] (a.south) -- (b.north);

      \draw[arr] (b.south) -- (c.north);

      \draw[arr] (c.south) -- (d.north);

    \end{tikzpicture}

  \end{center}

  \bigbreak

  Всё \enquote{до \texttt{project()}} --- в toolchain-файл или на
  командную строку.

\end{frame}

\begin{frame}[fragile]{Живое демо: \texttt{demo-toolchain.cmake}}

  \enquote{Кросс} на ту же архитектуру (Linux x86\_64
  $\rightarrow$ Linux x86\_64): доказательство механизма, а не
  архитектуры. Файл задаёт \texttt{CMAKE\_C\_COMPILER} и маркер
  \texttt{TOOLCHAIN\_MARKER}.

  \vspace{0.5em}

  \myinputlisting[minted language=cmake]
    {presentations/02-cmake/src/toolchain/}
    {demo-toolchain.cmake}

\end{frame}

\begin{frame}[fragile]{Живое демо: \texttt{demo-toolchain.cmake}}

  \myinputlisting[minted language=cmake]
    {presentations/02-cmake/src/}
    {calc/06/CMakeLists.txt}

\end{frame}

\begin{frame}[fragile]{Живое демо: \texttt{demo-toolchain.cmake}}

  \begin{terminalwindow}
!\shellcommand{cmake -S . -B build-tc \ \\ -DCMAKE\_TOOLCHAIN\_FILE=\dmtc}!
-- toolchain file read (before project())
-- toolchain file read (before project())
-- The C compiler identification is GNU 15.3.0
!\dots!
-- toolchain marker: loaded: !\dots!/demo-toolchain.cmake
-- CROSSCOMPILING = TRUE
  \end{terminalwindow}

\end{frame}

\begin{frame}[fragile]{\texttt{--toolchain}: вторая форма, порядок
  чтения}

  Вторая форма (CMake $\geq$ 3.21):
  \texttt{cmake \dots --toolchain <файл>} --- эквивалент
  \texttt{-DCMAKE\_TOOLCHAIN\_FILE=\dots}; результат тот же.

  \vspace{0.5em}

  \myinputlisting[minted language=cmake]
    {presentations/02-cmake/src/toolchain/order-demo/}
    {CMakeLists.txt}

  \vspace{0.5em}

  \begin{terminalwindow}
!\shellcommand{cmake -S . -B build-order \ \\ -DCMAKE\_TOOLCHAIN\_FILE=\dmtc}!
-- operator ABOVE project()
-- toolchain file read (before project())
-- toolchain file read (before project())
-- The C compiler identification is GNU 15.3.0
-- after project(): /usr/bin/gcc
  \end{terminalwindow}

\end{frame}

\begin{frame}[fragile]{Кросс-компиляция aarch64 (иллюстративно):
  toolchain-файл}

  Настоящий cross-toolchain: Linux x86\_64 $\rightarrow$ Linux aarch64.

  \textbf{Пример иллюстративный:} кросс-компилятора
  (\texttt{/usr/bin/aarch64-linux-gnu-gcc}) в окружении нет.
  Проверено лишь, что CMake \textbf{загружает} файл и
  \textbf{подхватывает} из него \texttt{CMAKE\_C\_COMPILER} (реальная
  ошибка ниже --- само доказательство); перекрестная сборка и вывод
  \texttt{file} \textbf{НЕ выполнялись}.

  \vspace{0.5em}

  \myinputlisting[minted language=cmake]
    {presentations/02-cmake/src/toolchain/}
    {aarch64-fragment.cmake}

\end{frame}

\begin{frame}[fragile]{Кросс-компиляция aarch64 (иллюстративно):
  результат}

  \begin{terminalwindow}
!\shellcommand{cmake -S . -B build-arm \ \\ -DCMAKE\_TOOLCHAIN\_FILE=\armtc}!
-- The C compiler identification is unknown
CMake Error at CMakeLists.txt:2 (project):
  The CMAKE_C_COMPILER:
    /usr/bin/aarch64-linux-gnu-gcc
  is not a full path to an existing compiler tool.
  \end{terminalwindow}

  \bigbreak

  Ожидаемый результат (\textbf{не выполнялся} --- иллюстрация того,
  что было бы с установленным кросс-компилятором и sysroot):

  \begin{terminalwindow}
!\shellcommand{cmake -S . -B build-arm \ \\ -DCMAKE\_TOOLCHAIN\_FILE=\armtc}!
!\shellcommand{cmake --build build-arm}!
!\shellcommand{file build-arm/calc\_app}!
build-arm/calc_app: ELF 64-bit LSB executable, ARM aarch64, !\dots!
  \end{terminalwindow}

  \bigbreak

  \underline{Итог блока: тот же \texttt{CMakeLists.txt}, та же
  \texttt{cmake --build} ---}

  \underline{\enquote{кросс} задаётся одним файлом
  на этапе configure.}

\end{frame}

\begin{frame}[fragile]{Подводный камень: \texttt{CMAKE\_C\_COMPILER}
  задан поздно}

  \begin{columns}[T]

    \begin{column}{0.5\textwidth}

      \begin{itemize}

        \item Подводный камень:
          \texttt{set(CMAKE\_C\_COMPILER \dots)} после
          \texttt{project()} --- поздно: компилятор уже определён
          (toolchain-файл читается при первом \texttt{project()}) и его
          определение кэшируется.

      \end{itemize}

      \myinputlisting[minted language=cmake,
        minted options={fontsize=\footnotesize, breaklines,
          breakanywhere, linenos, escapeinside=||, numbersep=2pt}]
        {presentations/02-cmake/src/toolchain/late-demo/}
        {CMakeLists.txt}

    \end{column}

    \begin{column}{0.48\textwidth}

      \begin{terminalwindow}
!\shellcommand{cmake -S . -B build-late}!
-- The C compiler identification is Clang 21.1.8
-- after set(): /usr/bin/gcc
      \end{terminalwindow}

      \vspace{0.4em}

      Clang выбран (хостовый \texttt{cc}); переменная после
      \texttt{project()} поменялась --- но сборка пойдёт на уже
      определённом компиляторе.

      \vspace{0.4em}

      Переконфигурация не меняет уже прочитанный toolchain-файл ---
      нужен новый каталог сборки.

    \end{column}

  \end{columns}

  \bigbreak

  \begin{center}

    \fbox{\textbf{Правило:} всё \enquote{до \texttt{project()}} ---
    в toolchain-файл или в \texttt{-D} на командной строке.}

  \end{center}

\end{frame}

\begin{frame}[fragile]{\texttt{install}: запуск}

  \begin{terminalwindow}
!\shellcommand{cmake --install build --prefix build/install}!
-- Install configuration: "Release"
-- Installing: !\dots!/build/install/bin/calc_app
-- Installing: !\dots!/build/install/lib/libcalc.a
-- Installing: !\dots!/build/install/include/calc.h
  \end{terminalwindow}

  \bigbreak

  \begin{terminalwindow}
!\shellcommand{find build/install -type f | sort}!
build/install/bin/calc_app
build/install/include/calc.h
build/install/lib/libcalc.a
!\shellcommand{./build/install/bin/calc\_app}!
!\myhighlight{yellow}{2 + 3 = 5}!
  \end{terminalwindow}

  \bigbreak

  Работает из install-места: статическая библиотека \enquote{вшита}
  в \texttt{calc\_app}.

\end{frame}

\begin{frame}{Шпаргалка: команды лекции}

  \begin{table}[h]

    \begin{center}

      \small

      \renewcommand{\arraystretch}{1.15}

      \begin{tabular}{|p{3.0cm}|p{9.6cm}|}

        \hline

        \textbf{Этап} & \textbf{Команда} \\

        \hline

        configure
          & \texttt{cmake -S . -B build
            -DCMAKE\_BUILD\_TYPE=Release} \\

        \hline

        configure с опциями
          & \texttt{\dots\ -DBUILD\_SHARED=ON} \\

        \hline

        configure, кросс
          & \texttt{\dots\ -DCMAKE\_TOOLCHAIN\_FILE=tc.cmake} \\

        \hline

        configure, кросс (форма 2)
          & \texttt{\dots\ --toolchain tc.cmake} \\

        \hline

        генератор
          & \texttt{cmake -S . -B build-ninja -G Ninja} \\

        \hline

        посмотреть опции
          & \texttt{cmake -L build} \\

        \hline

        переконфигурация
          & тот же \texttt{cmake -S . -B build}
            с новыми \texttt{-D\dots} \\

        \hline

        сборка
          & \texttt{cmake --build build} \\

        \hline

        сборка одной цели
          & \texttt{cmake --build build --target calc} \\

        \hline

        установка
          & \texttt{cmake --install build
            --prefix build/install} \\

        \hline

      \end{tabular}

    \end{center}

  \end{table}

  \vspace{0.4em}

  \small

  Одна пара команд на этап --- при любом генераторе и платформе:
  в этом и есть смысл мета-системы.

\end{frame}

\begin{frame}{Чего не было (следующие шаги)}

  \small

  \begin{itemize}

    \item \texttt{find\_package()} + config-пакеты
      (CPM/FetchContent) --- сторонние библиотеки: наш
      \texttt{calc} живёт в одном каталоге, реальные проекты
      тянут зависимости.

    \item \texttt{CTest} --- тесты, встроенные в сборку
      (\texttt{enable\_testing()}, \texttt{add\_test()}).

    \item \texttt{add\_custom\_command()} /
      \texttt{add\_custom\_target()} --- шаги, которые не
      являются компиляцией: генерация кода, упаковка, отчёты.

    \item \texttt{CMakePresets.json} --- именованные пресеты
      configure (генератор, опции, toolchain) в JSON: меньше
      копировать командную строку.

    \item \texttt{CPack} --- готовые пакеты
      (.deb/.rpm/\dots) из той же конфигурации.

  \end{itemize}

  \vspace{0.6em}

  По одной строке на тему: естественное продолжение, когда
  проект перерастает \texttt{add\_executable} +
  \texttt{target\_link\_libraries}.

\end{frame}

\begin{frame}{Итоги}

  \small

  \begin{columns}[T]

    \begin{column}{0.56\textwidth}

      \begin{itemize}

        \item CMake --- мета-система\footnotemark{} сборки:
          \texttt{CMakeLists.txt} (цели) $\rightarrow$
          генератор $\rightarrow$ \texttt{Makefile}/
          \texttt{build.ninja} $\rightarrow$ make/ninja
          (лекция 1 остаётся в силе, а не выбывается).

        \item Цикл: configure
          (\texttt{cmake -S . -B build}) $\rightarrow$ build
          (\texttt{cmake --build}) $\rightarrow$ install
          (\texttt{cmake --install});
          \texttt{CMakeCache.txt} --- память решений;
          reconfigure --- тот же \texttt{cmake}
          с новыми \texttt{-D}.

        \item Цели + \texttt{PUBLIC}/\texttt{PRIVATE}/
          \texttt{INTERFACE}: \enquote{требования к
          использованию} (включённые каталоги, библиотеки)
          переносятся по графу зависимостей --- писать
          \texttt{-I}/\texttt{-l} в клиентах не нужно.

        \item Переменные: обычные (область каталога) vs
          кэш-переменные (\texttt{-D} из командной строки);
          \texttt{option()} + \texttt{if()} ---
          настраиваемый проект.

        \item Toolchain-файл: читается при первом
          \texttt{project()}, до компиляторов, --- механизм
          кросс-компиляции (пример aarch64 ---
          иллюстративный).

        \item \texttt{calc/01/CMakeLists.txt} на 8 строк --- и
          кроссплатформенность в комплекте.

      \end{itemize}

      \footnotetext{\url{\cmakehelp/manual/cmake.1.html}}

    \end{column}

    \begin{column}{0.42\textwidth}

      \begin{center}

        \begin{tikzpicture}[
          box/.style={draw, thick, align=left, inner sep=2pt,
            font=\footnotesize, text width=3.5cm}
        ]

          \node[box] (s1) at (0,0)
            {\texttt{calc/01/CMakeLists.txt}\\8 строк};

          \node[box] (s2) at (0.35,1)
            {\texttt{calc/02/CMakeLists.txt}\\стандарт + флаги};

          \node[box] (s3) at (0.7,2)
            {\texttt{calc/03/CMakeLists.txt}\\option/if};

          \node[box] (s4) at (1.05,3)
            {\texttt{calc/04/CMakeLists.txt}\\install()};

          \node[box] (s5) at (1.4,4)
            {\texttt{calc/05/CMakeLists.txt}\\SHARED};

          \node[box] (s6) at (1.75,5)
            {\texttt{calc/06/CMakeLists.txt}\\+ toolchain};

          \draw[->, thick] (s1.north) -- (s2.south);

          \draw[->, thick] (s2.north) -- (s3.south);

          \draw[->, thick] (s3.north) -- (s4.south);

          \draw[->, thick] (s4.north) -- (s5.south);

          \draw[->, thick] (s5.north) -- (s6.south);

        \end{tikzpicture}

      \end{center}

    \end{column}

  \end{columns}

\end{frame}

\begin{frame}{Ссылки}

  \begin{itemize}

    \item Документация CMake (latest): \url{\cmakehelp/} --- основные
      узлы: \texttt{cmake(1)}, \texttt{cmake-language(7)},
      \texttt{cmake-buildsystem(7)}, \texttt{cmake-toolchains(7)},
      \texttt{cmake-generators(7)}.

    \item Примеры лекции: каталог \texttt{presentations/02-cmake/src/}
      презентации.

  \end{itemize}

\end{frame}
\end{document}
